% ============================================================
% Sección insertable en el Avance 2 (Seminario de Tesis Doctoral II)
% Procedimiento de búsqueda y selección de fuentes con el
% Agente de Fichas Bibliográficas v11
%
% Requiere en el preámbulo:
%   \usepackage{booktabs,longtable,array,url}
% La tabla de bitácora se incluye con:  % Generado por scripts/generar_bitacora.py — no editar a mano.
% Requiere: \usepackage{longtable,booktabs,array}
\begin{longtable}{@{}p{0.55cm}p{3.2cm}p{3.6cm}p{1.9cm}p{1.1cm}p{0.6cm}p{3.1cm}@{}}
\caption{Bitácora de fichas aprobadas por el Agente de Fichas Bibliográficas v11 (decisión humana registrada).}\label{tab:bitacora-corpus}\\
\toprule Dim. & Fuente & DOI & Verificación & Citas & AA & Comentario de decisión \\ \midrule \endfirsthead
\toprule Dim. & Fuente & DOI & Verificación & Citas & AA & Comentario de decisión \\ \midrule \endhead
\bottomrule \endfoot
D1 & Kasneci et al. (2023) & {\scriptsize 10.1016/j.lindif.2023.102274} & VERIFICADA & 4511 & Sí & {\small ética de LLMs en educación} \\
D1 & Li et al. (2025) & {\scriptsize 10.1371/journal.pone.0333411} & VERIFICADA & 11 & Sí & {\small aporta marco de transparencia algorítmica'} \\
D1 & Zawacki-Richter et al. (2019) & {\scriptsize 10.1186/s41239-019-0171-0} & VERIFICADA & 4714 & Sí & {\small NO HAY IA ENN LA DOCENCIA} \\
D1 & Chan et al. (2023) & {\scriptsize 10.1186/s41239-023-00411-8} & VERIFICADA & 1573 & Sí & {\small ok} \\
D1 & Dwivedi et al. (2023) & {\scriptsize 10.1016/j.ijinfomgt.2023.102642} & VERIFICADA & 3421 & Sí & {\small (URL de acceso registrada)} \\
D1 & Johnston et al. (2024) & {\scriptsize 10.1007/s40979-024-00149-4} & VERIFICADA & 219 & Sí & {\small a} \\
D1 & Acosta Enríquez et al. (2024) & {\scriptsize 10.1007/s40979-024-00157-4} & VERIFICADA & 82 & Sí & {\small arco etico cultural} \\
D1 & Plata et al. (2023) & {\scriptsize 10.24191/ajue.v19i4.24697} & VERIFICADA & 46 & Sí & {\small tendencias de investigación en integridad académica} \\
D1 & Wang et al. (2024) & {\scriptsize 10.1016/j.caeai.2024.100326} & VERIFICADA & 130 & Sí & {\small JUSTO LO QUE BUCAMOS} \\
D1 & Shah et al. (2025) & {\scriptsize Verificar en fuente} & NO ENCONTRADA & 0 & No & {\small (URL de acceso registrada)} \\
D1 & Watson et al. (2025) & {\scriptsize 10.1007/s00146-024-02174-w} & VERIFICADA & 22 & Sí & {\small porta marco de transparencia algorítmica'} \\
D2 & Barredo Arrieta et al. (2019) & {\scriptsize 10.1016/j.inffus.2019.12.012} & VERIFICADA & 8490 & Sí & {\small DA COPNCEPTOS DE MACHINE LEARNING} \\
D2 & Rughiniș et al. (2025) & {\scriptsize 10.3389/fcomp.2025.1608276} & VERIFICADA & 3 & Sí & {\small ave para D2 — aporta marco de transparencia algorítmica} \\
D2 & Ali et al. (2023) & {\scriptsize 10.1016/j.inffus.2023.101805} & VERIFICADA & 1370 & Sí & {\small Discussion of different explainability assessment methods.} \\
D2 & Floridi et al. (2018) & {\scriptsize 10.1007/s11023-018-9482-5} & VERIFICADA & 3156 & Sí & {\small marco etico} \\
D2 & Seeger et al. (2026) & {\scriptsize 10.1007/s12525-026-00883-2} & VERIFICADA & 0 & Sí & {\small APORTA UN EJEMPLO DE ETIQUETADO  EN SOCUIAL MEDIA} \\
D2 & Hagendorff (2020) & {\scriptsize 10.1007/s11023-020-09517-8} & VERIFICADA & 1558 & Sí & {\small guia de marco etico} \\
D2 & Mohamed (2024) & {\scriptsize 10.29140/tltl.v6n3.1813} & VERIFICADA & 27 & Sí & {\small TRANSPARENCIA} \\
D2 & Arriaga Navarro (2026) & {\scriptsize --} & SIN DOI — URL manual & 0 & No & {\small aporta marco de transparencia algorítmica} \\
D2 & Enríquez Jiménez (2024) & {\scriptsize https://hdl.handle.net/20.500.12371/24132} & NO ENCONTRADA & 0 & No & {\small porta marco de transparencia algorítmica'} \\
D2 & Verificar en fuente (2023) & {\scriptsize 10.1186/s41239-023-00408-3} & VERIFICADA & 1082 & Sí & {\small importante inferencia en todo} \\
D3 & Cerro Martínez et al. (2020) & {\scriptsize 10.1186/s41239-020-00217-y} & VERIFICADA & 45 & Sí & {\small analitica} \\
D3 & Grové (2021) & {\scriptsize 10.3389/fpsyt.2020.606041} & VERIFICADA & 129 & Sí & {\small CODISEÑO} \\
D3 & Verificar en fuente completa para autores y detalles. & {\scriptsize 10.1111/ejed.12533} & VERIFICADA & 671 & Sí & {\small d2} \\
D3 & Axelsson et al. (2021) & {\scriptsize 10.1145/3472225} & VERIFICADA & 79 & Sí & {\small FRAMEWORK DE DISEÑO TECNOLOGICO} \\
D3 & Nicholson et al. (2022) & {\scriptsize 10.1145/3491102.3517667} & VERIFICADA & 34 & No & {\small CO ENSEÑANZA} \\
D3 & Romero et al. (2020) & {\scriptsize 10.3390/su12187798} & VERIFICADA & 82 & Sí & {\small sistema de etiquetado indagacion} \\
D3 & Engelbertink et al. (2020) & {\scriptsize 10.1007/s10639-020-10147-x} & VERIFICADA & 39 & Sí & {\small PODEMOS SACAR UN CASOP DE COMO AYUDA  EL BLEND LEARNIG YPARA QUE NLOS } \\
D3 & Smith et al. (2022) & {\scriptsize 10.1007/s10639-021-10845-0} & VERIFICADA & 16 & Sí & {\small DISEÑO} \\
D3 & Carcamo Garcia (2024) & {\scriptsize --} & SIN DOI — URL manual & 0 & No & {\small USO ACADEMICO} \\
D3 & Moeini (2020) & {\scriptsize 10.13140/rg.2.2.32813.61921} & VERIFICADA & 3 & Sí & {\small diseño de etiquetado} \\
D3 & Parker et al. (2025) & {\scriptsize 10.1007/s10791-025-09745-5} & VERIFICADA & 3 & Sí & {\small ave para D2 — aporta marco de transparencia algorítmica'} \\
D4 & Ponce Rojo et al. (2025) & {\scriptsize 10.3390/educsci15091134} & VERIFICADA & 4 & Sí & {\small aporta marco de transparencia algorítmica} \\
D4 & Rivadeneira et al. (2024) & {\scriptsize 10.1145/3689370} & VERIFICADA & 8 & Sí & {\small LOS CRITERIOS BAJO LOS QUE USAN CHAT GPT} \\
D4 & Juarez Merino (2026) & {\scriptsize 10.18848/1832-3669/cgp/a1024} & VERIFICADA & 0 & Sí & {\small Análisis crítico de IA como reconfigurador de poder y gobernanza en co} \\
D4 & Miralrio et al. (2024) & {\scriptsize 10.3389/fcomp.2024.1250992} & VERIFICADA & 8 & Sí & {\small Gran potencial educativo pero debate continuo sobre riesgos sociales.} \\
D4 & Durso et al. (2022) & {\scriptsize 10.33225/pec/22.80.679} & VERIFICADA & 39 & Sí & {\small Brasil tiene mayor mercado de educacion superior en Latinoamerica} \\
D4 & Haroud et al. (2025) & {\scriptsize 10.3390/educsci15040396} & VERIFICADA & 43 & Sí & {\small ebates sobre potencial transformador de practicas pedagogicas y redefi} \\
D4 & Aguilar Rivera (2025) & {\scriptsize --} & SIN DOI — URL manual & 0 & No & {\small Concepto VIPOGIA y marco regulador basado en principios derechos human} \\
D4 & Henrickson (2021) & {\scriptsize 10.26174/thesis.lboro.10596404.v1} & VERIFICADA & 15 & Sí & {\small Marco hermeneutico para comprender autoria e impacto social de conteni} \\
D4 & Verificar en fuente (Verificar en fuente) & {\scriptsize --} & SIN DOI - fuente manual & 0 & No & {\small aporta marco de transparencia algorítmica} \\
\end{longtable}

\begin{table}[htbp]\centering\small
\caption{Resumen cuantitativo del corpus generado con el agente.}\label{tab:resumen-corpus}
\begin{tabular}{@{}lr@{}}\toprule Indicador & Valor \\ \midrule
Fichas aprobadas & 41 \\
\quad D1 · Gobernanza y ética de la IA en educación & 11 \\
\quad D2 · Transparencia y etiquetado de IA en contenidos & 10 \\
\quad D3 · Diseño participativo en tecnología educativa & 11 \\
\quad D4 · Perspectiva crítica IA educación latinoamericana & 9 \\
Fuentes rechazadas con DOI & 19 \\
Fuentes rechazadas registradas por título & 23 \\
DOI verificado & 35 \\
DOI no localizado & 2 \\
Sin DOI (tesis de repositorio, modo manual) & 4 \\
Acceso abierto & 34 \\
Fechas de sesión & 19/04/2026, 22/04/2026 \\
\bottomrule\end{tabular}\end{table}

% (archivo generado por scripts/generar_bitacora.py)
% ============================================================

\section{Procedimiento de búsqueda y selección de fuentes}
\label{sec:agente}

Atendiendo a la observación de que la estrategia del Avance~1 no constituía
una búsqueda reproducible, esta sección describe la herramienta con la que se
construyó el corpus del estado del arte, el procedimiento que ejecuta y las
evidencias que deja. La descripción sigue los campos de reporte por base que
recomienda la extensión PRISMA-S (Rethlefsen et al., 2021): ecuación literal,
campos consultados, filtros, fecha, registros recuperados, duplicados,
cribado y colección final.

\subsection{Identificación de la herramienta}

La búsqueda se ejecutó con el \emph{Agente de Fichas Bibliográficas},
versión~11, un programa propio escrito en Node.js y versionado en Git, de
modo que cada modificación del procedimiento queda fechada. El agente consulta
cuatro bases abiertas, sin inicio de sesión ni extracción de páginas web, y
emplea el modelo de lenguaje Claude Sonnet~4 (Anthropic, identificador
\texttt{claude-sonnet-4-20250514}) únicamente para redactar la ficha
estructurada a partir de los metadatos recuperados. El modelo no busca ni
decide la inclusión. El código del agente se desarrolló con apoyo de Claude
Code (Anthropic), lo que consta en el historial del repositorio.

\begin{table}[htbp]\centering\small
\caption{Bases consultadas por el agente y su naturaleza.}
\label{tab:bases}
\begin{tabular}{@{}p{3.1cm}p{5.3cm}p{5.6cm}@{}}
\toprule
Base & Naturaleza & Uso en el agente \\
\midrule
DOAJ & Directorio de revistas arbitradas de acceso abierto & Primera consulta para artículos; garantiza texto completo disponible \\
OpenAlex & Índice bibliográfico abierto & Artículos y tesis; verificación de DOI, citas y acceso abierto \\
Crossref & Agencia de registro de DOI con metadatos editoriales & Artículos de revista; verificación de DOI de respaldo \\
Semantic Scholar & Grafo académico con resúmenes e identificadores & Artículos y tesis; citas y PDF de acceso abierto \\
\bottomrule
\end{tabular}
\end{table}

Esta distinción corrige el Cuadro~12 del Avance~1: Springer Nature, Elsevier,
MDPI y PLoS son las editoriales de los artículos recuperados, no las bases en
las que se ejecutó la búsqueda.

\subsection{Lógica del procedimiento}

Cada sesión del agente ejecuta cinco etapas; las cuatro primeras son
automáticas y la quinta es siempre humana:
(1)~recuperación en las cuatro bases con la ecuación y la página que
corresponden a la sesión;
(2)~depuración: deduplicación por DOI y por título, exclusión de fuentes ya
aprobadas o rechazadas, orden por citas y corte en veinte candidatos;
(3)~redacción de $N$ fichas ($2 \le N \le 5$) por el modelo de lenguaje a
partir de los metadatos reales;
(4)~verificación del DOI de cada ficha en OpenAlex y, en su defecto, en
Crossref, con registro de estado, título real, citas y acceso abierto;
(5)~decisión de la autora sobre cada ficha (aprobar, rechazar de forma
permanente o saltar) con comentario, fecha y hora.

Tres propiedades sostienen la trazabilidad. Ninguna fuente entra al corpus
sin decisión humana registrada. Ninguna fuente se evalúa dos veces, porque las
aprobadas y las rechazadas se excluyen de las búsquedas posteriores por DOI o
por título. La búsqueda avanza de forma determinista: el archivo de sesiones
guarda, por dimensión y tipo documental, la ecuación y la página consultadas,
y cada ecuación se recorre tres páginas antes de rotar a la siguiente.

\subsection{Especificación reproducible}

El corpus se organiza en las cuatro dimensiones definidas en Seminario de
Tesis~I. Para cada una, el agente tiene seis ecuaciones en inglés para
artículos y cuatro en español para tesis, que se transcriben en el
Anexo~\ref{anx:ecuaciones}. Son cadenas de términos en lenguaje natural, no
expresiones booleanas, porque las cuatro bases operan con motores de
relevancia; esto sustituye la fórmula booleana del Anexo~C del Avance~1, que
arrojaba más de catorce mil registros en Scopus.

\begin{table}[htbp]\centering\scriptsize
\caption{Sintaxis, campos y filtros por base.}
\label{tab:sintaxis}
\begin{tabular}{@{}p{2.2cm}p{3.2cm}p{3.0cm}p{3.3cm}p{1.6cm}p{0.9cm}@{}}
\toprule
Base & Punto de acceso & Campo consultado & Filtros fijos & Orden & Reg./pág. \\
\midrule
DOAJ & \texttt{/api/v3/search/articles/} & Registro completo & Ninguno & Relevancia & 8 \\
OpenAlex (artículos) & \texttt{/works?search=} & Título, resumen, texto indexado & \texttt{publication\_year:2018-2025}, \texttt{type:article} & Relevancia & 10 \\
OpenAlex (tesis) & \texttt{/works?search=} & Ídem & \texttt{publication\_year:2018-2025}, \texttt{type:dissertation} & Relevancia & 10 \\
Crossref & \texttt{/works?query=} & Metadatos bibliográficos & \texttt{from-pub-date:2018}, \texttt{type:journal-article} & Relevancia & 8 \\
Semantic Scholar & \texttt{/graph/v1/paper/search} & Título y resumen & Ninguno; para tesis se añade \emph{thesis dissertation} y se filtra por tipo & Relevancia & 8 \\
\bottomrule
\end{tabular}
\end{table}

De cada registro se recuperan título, hasta tres autores, revista o
institución, año, resumen truncado a 500 caracteres, DOI, citas y base de
origen. Dos registros se consideran duplicados si comparten DOI o los
primeros cincuenta caracteres del título; prevalece el recuperado primero, por
lo que se conserva la versión de acceso abierto cuando existe. Una sesión
completa recupera como máximo 52 registros antes de deduplicar y presenta 20
al modelo.

La instrucción de sistema que recibe el modelo impone seis reglas: usar
exclusivamente los datos entregados; extraer los hallazgos del resumen y, si
no lo hay, escribir «Sin resumen, verificar en fuente»; no generar URL;
priorizar 2020--2025 y mayor número de citas; limitar cada campo a sesenta
palabras; y devolver sólo JSON. La instrucción de usuario indica la
dimensión, el número de fichas y la lista de candidatos, y pide seleccionar
las $N$ más relevantes para el proyecto. Las fichas contienen folio, tipo,
referencia, DOI, dimensión, revista y base de origen, contexto, año,
paradigma, metodología, objeto de estudio, tres hallazgos, aportación,
brecha, pertinencia para el marco teórico, justificación de inclusión y
resumen preliminar.

Las tesis de TESIUNAM, BUAP o UAEM, que no tienen DOI en estas bases, se
incorporan por un modo manual en el que la autora pega el registro del
repositorio; la ficha queda marcada con origen «Manual» y estado «sin DOI».

\subsection{Evolución de las versiones}

El procedimiento se calibró a lo largo de varias versiones. La versión~6
consultaba OpenAlex y Crossref con un término por dimensión y generaba tres
fichas por sesión; su limitación era que cada sesión repetía los primeros
resultados y no conservaba memoria de lo ya evaluado. La versión~11 incorporó
la paginación automática, la rotación de ecuaciones, la lista de rechazadas
permanentes, el modo manual para tesis de repositorios mexicanos y Semantic
Scholar como tercera base. Sobre esa línea base, el historial del
repositorio fecha el 26 de octubre de 2026 cuatro cambios: la interfaz web
para mostrar el corpus y las decisiones; la incorporación de DOAJ como
cuarta base y, acto seguido, como primera en el orden de consulta, para que
en los duplicados prevalezca la versión de acceso abierto; y los campos
obligatorios de justificación de inclusión y resumen preliminar en cada
ficha. [Las versiones intermedias 7 a 10 se completan con los registros de
la autora.]

\subsection{Discriminación automática y embudo de selección}

Las fuentes rechazadas por la autora son la última etapa de un embudo cuya
mayor parte ejecuta el agente por reglas declaradas. Los filtros codificados
son: ventana temporal desde 2018 (con prioridad 2020--2025 en la instrucción
al modelo); tipo documental restringido a artículos de revista y tesis, lo
que excluye blogs, notas de prensa y materiales sin respaldo metodológico;
preferencia por acceso abierto mediante la consulta prioritaria a DOAJ;
identificador verificable, pues toda fuente automática debe tener DOI y las
que carecen de él sólo entran por modo manual; exclusión de lo ya revisado;
corte por relevancia (tres páginas por ecuación), orden por citas y
presentación de veinte candidatos; y selección por el modelo de sólo $N$
fichas por sesión. Además, en la revisión la autora aplicó tres criterios de
descarte rápido que no están codificados como exclusión dura: ausencia de
DOI, ausencia de versión descargable y antigüedad mayor a tres años salvo
referencia obligada. El corpus muestra que operaron como preferencia y no
como filtro: contiene cuatro tesis sin DOI, siete fuentes sin versión de
acceso abierto y once fuentes de 2018--2021.

\begin{table}[htbp]\centering\small
\caption{Embudo de una sesión del agente (cotas fijadas por el código).}
\label{tab:embudo}
\begin{tabular}{@{}p{7.2cm}p{3.4cm}p{3.0cm}@{}}
\toprule
Etapa & Registros & Quién descarta \\
\midrule
Hits totales de la ecuación en las cuatro bases & Miles (no se visitan) & Corte de paginación \\
Recuperados en la página de la sesión & Hasta 52 & --- \\
Tras deduplicar y excluir ya revisadas & Variable & Agente \\
Presentados al modelo & 20 & Agente (orden por citas) \\
Fichas redactadas & 2 a 5 (por defecto 3) & Modelo \\
Aprobadas & 0 a 3 & Autora \\
\bottomrule
\end{tabular}
\end{table}

Con 64 decisiones registradas (41 aprobadas y 23 rechazadas) y tres fichas
por sesión, se estiman unas 21 sesiones, unos 420 candidatos presentados al
modelo y hasta unos 1\,100 registros brutos recuperados. Las 23 rechazadas a
mano representan menos del 6\,\% de lo descartado. Esta cifra es una
estimación porque la versión~11 no persiste los conteos por sesión; el
script \texttt{conteos\_por\_base.mjs} del repositorio obtiene el total de
registros por ecuación y base con los filtros del agente, y su salida se
anexará como «registros recuperados» conforme a PRISMA-S.

\subsection{Evidencias y colección final}

El agente deja cuatro archivos: el corpus aprobado (JSON y texto), la lista
de rechazadas, el estado de sesiones y el historial de cambios del
procedimiento. A partir de ellos se genera la bitácora de la
Tabla~\ref{tab:bitacora-corpus} (Anexo). El resumen cuantitativo se presenta
en la Tabla~\ref{tab:resumen-corpus}. Dos de las cinco fuentes de la matriz
inicial del artículo provienen del corpus del agente (Seeger et al., 2026,
folio D2-A02; Axelsson et al., 2021, folio D3-A02); las fuentes del
etiquetado alimentario se localizaron por revisión de listas de referencias y
constituyen el estrato analógico, que se reporta por separado.

\subsection{Límites del registro y ajustes derivados}

La versión~11 no persiste la fecha de cada consulta, el total de registros
que devuelve cada base ni el motivo de cada rechazo; tampoco lee el texto
completo, pues la ficha se redacta a partir del resumen. En consecuencia, la
estrategia debe declararse como muestreo por relevancia (tres primeras
páginas por ecuación) y no como búsqueda exhaustiva, lo que es congruente con
una revisión de alcance; los campos de diseño, muestra, instrumentos y
hallazgos de la matriz se extraen de la lectura completa y no de la ficha
automática. Dos ajustes se derivan del análisis del corpus. Primero, las
ecuaciones de D2 contenían «transparency» y «disclosure» y atrajeron
literatura de explicabilidad; sólo Seeger et al. (2026) estudia una etiqueta
de contenido generado con IA. La afirmación del Avance~1 sobre un único
antecedente directo se retira: las referencias señaladas por el asesor
muestran que la literatura existe y que las ecuaciones no incluían su
vocabulario. Segundo, la ventana se ajusta a 2022--2026 para la evidencia
directa, con 2018--2021 como antecedentes, y se incorpora un bloque de
ecuaciones sobre diseño y evaluación de etiquetas de contenido sintético.

% ------------------------------------------------------------
\appendix
\section{Ecuaciones de búsqueda del agente}
\label{anx:ecuaciones}

\begin{description}\small
\item[D1 Gobernanza y ética de la IA en educación.]
Artículos: \emph{AI governance higher education ethics transparency};
\emph{artificial intelligence academic integrity university policy};
\emph{responsible AI education institutional governance};
\emph{AI ethics higher education accountability};
\emph{algorithmic accountability academic institutions policy};
\emph{generative AI governance university Mexico}.
Tesis: \emph{gobernanza inteligencia artificial educacion superior};
\emph{etica inteligencia artificial universidad Mexico};
\emph{uso responsable IA educacion superior posgrado};
\emph{politicas institucionales inteligencia artificial universidad}.

\item[D2 Transparencia y etiquetado de IA en contenidos.]
Artículos: \emph{AI disclosure labeling academic content transparency};
\emph{artificial intelligence transparency labeling education};
\emph{algorithmic transparency content authorship higher education};
\emph{AI generated content labeling standards};
\emph{machine learning content disclosure academic writing};
\emph{generative AI authorship transparency university}.
Tesis: \emph{transparencia inteligencia artificial contenidos educativos};
\emph{etiquetado IA materiales educativos universidad};
\emph{declaracion autoria inteligencia artificial academia Mexico};
\emph{sistema etiquetado contenidos generados IA educacion}.

\item[D3 Diseño participativo en tecnología educativa.]
Artículos: \emph{participatory design educational technology usability};
\emph{co-design learning materials higher education};
\emph{heuristic evaluation usability educational materials Nielsen};
\emph{user centered design learning tools students};
\emph{participatory action research educational technology};
\emph{co-creation digital learning resources university}.
Tesis: \emph{diseno participativo tecnologia educativa universidad};
\emph{co-diseno materiales educativos digitales posgrado};
\emph{usabilidad recursos educativos digitales Mexico};
\emph{evaluacion heuristica materiales aprendizaje}.

\item[D4 Perspectiva crítica de la IA en la educación latinoamericana.]
Artículos: \emph{critical educational technology AI Latin America};
\emph{critical digital literacy higher education AI};
\emph{AI education inequality critical perspective Mexico};
\emph{technology education power critical pedagogy};
\emph{digital colonialism education artificial intelligence};
\emph{critical algorithm studies education equity}.
Tesis: \emph{tecnologia educativa critica America Latina inteligencia artificial};
\emph{pedagogia critica inteligencia artificial Mexico universidad};
\emph{alfabetizacion digital critica educacion superior};
\emph{desigualdad tecnologica educacion superior Mexico}.
\end{description}

\section{Bitácora del corpus}
% Generado por scripts/generar_bitacora.py — no editar a mano.
% Requiere: \usepackage{longtable,booktabs,array}
\begin{longtable}{@{}p{0.55cm}p{3.2cm}p{3.6cm}p{1.9cm}p{1.1cm}p{0.6cm}p{3.1cm}@{}}
\caption{Bitácora de fichas aprobadas por el Agente de Fichas Bibliográficas v11 (decisión humana registrada).}\label{tab:bitacora-corpus}\\
\toprule Dim. & Fuente & DOI & Verificación & Citas & AA & Comentario de decisión \\ \midrule \endfirsthead
\toprule Dim. & Fuente & DOI & Verificación & Citas & AA & Comentario de decisión \\ \midrule \endhead
\bottomrule \endfoot
D1 & Kasneci et al. (2023) & {\scriptsize 10.1016/j.lindif.2023.102274} & VERIFICADA & 4511 & Sí & {\small ética de LLMs en educación} \\
D1 & Li et al. (2025) & {\scriptsize 10.1371/journal.pone.0333411} & VERIFICADA & 11 & Sí & {\small aporta marco de transparencia algorítmica'} \\
D1 & Zawacki-Richter et al. (2019) & {\scriptsize 10.1186/s41239-019-0171-0} & VERIFICADA & 4714 & Sí & {\small NO HAY IA ENN LA DOCENCIA} \\
D1 & Chan et al. (2023) & {\scriptsize 10.1186/s41239-023-00411-8} & VERIFICADA & 1573 & Sí & {\small ok} \\
D1 & Dwivedi et al. (2023) & {\scriptsize 10.1016/j.ijinfomgt.2023.102642} & VERIFICADA & 3421 & Sí & {\small (URL de acceso registrada)} \\
D1 & Johnston et al. (2024) & {\scriptsize 10.1007/s40979-024-00149-4} & VERIFICADA & 219 & Sí & {\small a} \\
D1 & Acosta Enríquez et al. (2024) & {\scriptsize 10.1007/s40979-024-00157-4} & VERIFICADA & 82 & Sí & {\small arco etico cultural} \\
D1 & Plata et al. (2023) & {\scriptsize 10.24191/ajue.v19i4.24697} & VERIFICADA & 46 & Sí & {\small tendencias de investigación en integridad académica} \\
D1 & Wang et al. (2024) & {\scriptsize 10.1016/j.caeai.2024.100326} & VERIFICADA & 130 & Sí & {\small JUSTO LO QUE BUCAMOS} \\
D1 & Shah et al. (2025) & {\scriptsize Verificar en fuente} & NO ENCONTRADA & 0 & No & {\small (URL de acceso registrada)} \\
D1 & Watson et al. (2025) & {\scriptsize 10.1007/s00146-024-02174-w} & VERIFICADA & 22 & Sí & {\small porta marco de transparencia algorítmica'} \\
D2 & Barredo Arrieta et al. (2019) & {\scriptsize 10.1016/j.inffus.2019.12.012} & VERIFICADA & 8490 & Sí & {\small DA COPNCEPTOS DE MACHINE LEARNING} \\
D2 & Rughiniș et al. (2025) & {\scriptsize 10.3389/fcomp.2025.1608276} & VERIFICADA & 3 & Sí & {\small ave para D2 — aporta marco de transparencia algorítmica} \\
D2 & Ali et al. (2023) & {\scriptsize 10.1016/j.inffus.2023.101805} & VERIFICADA & 1370 & Sí & {\small Discussion of different explainability assessment methods.} \\
D2 & Floridi et al. (2018) & {\scriptsize 10.1007/s11023-018-9482-5} & VERIFICADA & 3156 & Sí & {\small marco etico} \\
D2 & Seeger et al. (2026) & {\scriptsize 10.1007/s12525-026-00883-2} & VERIFICADA & 0 & Sí & {\small APORTA UN EJEMPLO DE ETIQUETADO  EN SOCUIAL MEDIA} \\
D2 & Hagendorff (2020) & {\scriptsize 10.1007/s11023-020-09517-8} & VERIFICADA & 1558 & Sí & {\small guia de marco etico} \\
D2 & Mohamed (2024) & {\scriptsize 10.29140/tltl.v6n3.1813} & VERIFICADA & 27 & Sí & {\small TRANSPARENCIA} \\
D2 & Arriaga Navarro (2026) & {\scriptsize --} & SIN DOI — URL manual & 0 & No & {\small aporta marco de transparencia algorítmica} \\
D2 & Enríquez Jiménez (2024) & {\scriptsize https://hdl.handle.net/20.500.12371/24132} & NO ENCONTRADA & 0 & No & {\small porta marco de transparencia algorítmica'} \\
D2 & Verificar en fuente (2023) & {\scriptsize 10.1186/s41239-023-00408-3} & VERIFICADA & 1082 & Sí & {\small importante inferencia en todo} \\
D3 & Cerro Martínez et al. (2020) & {\scriptsize 10.1186/s41239-020-00217-y} & VERIFICADA & 45 & Sí & {\small analitica} \\
D3 & Grové (2021) & {\scriptsize 10.3389/fpsyt.2020.606041} & VERIFICADA & 129 & Sí & {\small CODISEÑO} \\
D3 & Verificar en fuente completa para autores y detalles. & {\scriptsize 10.1111/ejed.12533} & VERIFICADA & 671 & Sí & {\small d2} \\
D3 & Axelsson et al. (2021) & {\scriptsize 10.1145/3472225} & VERIFICADA & 79 & Sí & {\small FRAMEWORK DE DISEÑO TECNOLOGICO} \\
D3 & Nicholson et al. (2022) & {\scriptsize 10.1145/3491102.3517667} & VERIFICADA & 34 & No & {\small CO ENSEÑANZA} \\
D3 & Romero et al. (2020) & {\scriptsize 10.3390/su12187798} & VERIFICADA & 82 & Sí & {\small sistema de etiquetado indagacion} \\
D3 & Engelbertink et al. (2020) & {\scriptsize 10.1007/s10639-020-10147-x} & VERIFICADA & 39 & Sí & {\small PODEMOS SACAR UN CASOP DE COMO AYUDA  EL BLEND LEARNIG YPARA QUE NLOS } \\
D3 & Smith et al. (2022) & {\scriptsize 10.1007/s10639-021-10845-0} & VERIFICADA & 16 & Sí & {\small DISEÑO} \\
D3 & Carcamo Garcia (2024) & {\scriptsize --} & SIN DOI — URL manual & 0 & No & {\small USO ACADEMICO} \\
D3 & Moeini (2020) & {\scriptsize 10.13140/rg.2.2.32813.61921} & VERIFICADA & 3 & Sí & {\small diseño de etiquetado} \\
D3 & Parker et al. (2025) & {\scriptsize 10.1007/s10791-025-09745-5} & VERIFICADA & 3 & Sí & {\small ave para D2 — aporta marco de transparencia algorítmica'} \\
D4 & Ponce Rojo et al. (2025) & {\scriptsize 10.3390/educsci15091134} & VERIFICADA & 4 & Sí & {\small aporta marco de transparencia algorítmica} \\
D4 & Rivadeneira et al. (2024) & {\scriptsize 10.1145/3689370} & VERIFICADA & 8 & Sí & {\small LOS CRITERIOS BAJO LOS QUE USAN CHAT GPT} \\
D4 & Juarez Merino (2026) & {\scriptsize 10.18848/1832-3669/cgp/a1024} & VERIFICADA & 0 & Sí & {\small Análisis crítico de IA como reconfigurador de poder y gobernanza en co} \\
D4 & Miralrio et al. (2024) & {\scriptsize 10.3389/fcomp.2024.1250992} & VERIFICADA & 8 & Sí & {\small Gran potencial educativo pero debate continuo sobre riesgos sociales.} \\
D4 & Durso et al. (2022) & {\scriptsize 10.33225/pec/22.80.679} & VERIFICADA & 39 & Sí & {\small Brasil tiene mayor mercado de educacion superior en Latinoamerica} \\
D4 & Haroud et al. (2025) & {\scriptsize 10.3390/educsci15040396} & VERIFICADA & 43 & Sí & {\small ebates sobre potencial transformador de practicas pedagogicas y redefi} \\
D4 & Aguilar Rivera (2025) & {\scriptsize --} & SIN DOI — URL manual & 0 & No & {\small Concepto VIPOGIA y marco regulador basado en principios derechos human} \\
D4 & Henrickson (2021) & {\scriptsize 10.26174/thesis.lboro.10596404.v1} & VERIFICADA & 15 & Sí & {\small Marco hermeneutico para comprender autoria e impacto social de conteni} \\
D4 & Verificar en fuente (Verificar en fuente) & {\scriptsize --} & SIN DOI - fuente manual & 0 & No & {\small aporta marco de transparencia algorítmica} \\
\end{longtable}

\begin{table}[htbp]\centering\small
\caption{Resumen cuantitativo del corpus generado con el agente.}\label{tab:resumen-corpus}
\begin{tabular}{@{}lr@{}}\toprule Indicador & Valor \\ \midrule
Fichas aprobadas & 41 \\
\quad D1 · Gobernanza y ética de la IA en educación & 11 \\
\quad D2 · Transparencia y etiquetado de IA en contenidos & 10 \\
\quad D3 · Diseño participativo en tecnología educativa & 11 \\
\quad D4 · Perspectiva crítica IA educación latinoamericana & 9 \\
Fuentes rechazadas con DOI & 19 \\
Fuentes rechazadas registradas por título & 23 \\
DOI verificado & 35 \\
DOI no localizado & 2 \\
Sin DOI (tesis de repositorio, modo manual) & 4 \\
Acceso abierto & 34 \\
Fechas de sesión & 19/04/2026, 22/04/2026 \\
\bottomrule\end{tabular}\end{table}

